\subsection{Implementacja}
\addcontentsline{toc}{subsection}{Implementacja}
\subsubsection{Scneariusz dochodów ze sprzedaży}
\addcontentsline{toc}{subsubsection}{Scneariusz dochodów ze sprzedaży}
Przychody ze sprzedaży poszczególnych typów produktów definiowane są przez wektor losowy opisany w treści zadania. W celu wygenerowania wektorów 
reprezentujących poszczególne scenariusze przychodów zastosowano bibliotekę MASS języka R. Implementacja została wykonana w środowisku R Studio IDE, a 
skrypt generujący dane zapisano w pliku \textit{t-student.R}. W ramach przeprowadzonej symulacji wygenerowano 
1000 scenariuszy realizacji przychodów.

\subsubsection{Model}
\addcontentsline{toc}{subsubsection}{Model}
Model zaimplementowano w środowisku IBM ILOG CPLEX Optimization Studio z wykorzystaniem solvera CPLEX. Nazewnictwo parametrów oraz zmiennych decyzyjnych 
jest zgodne z opisem zawartym w tabelach \ref{tab:param} i \ref{tab:var}. Plik \textit{wdwr25406-1.dat} zawiera definicje parametrów modelu, 
natomiast plik \textit{wdwr25406-1.mod} obejmuje wczytywanie parametrów, definicje zmiennych decyzyjnych, funkcji celu oraz ograniczeń modelu. 
W celu uproszczenia implementacji przyjęto numeryczne oznaczenia dla miesięcy, produktów oraz procesów technologicznych. Miesiące numerowane są chronologicznie, 
produkty zgodnie z indeksem występującym w nazwie (P1-P4), natomiast procesy technologiczne według poniższej sekwencji:
\begin{enumerate}
\item Szlifowanie,
\item Wiercenie pionowe,
\item Wiercenie poziome,
\item Frezowanie,
\item Toczenie.
\end{enumerate}

